\documentclass[msc]{VThesis}
%        اگر قصد نوشتن پروژه کارشناسی را دارید، در خط زیر به جای msc، کلمه bsc و اگر قصد نوشتن رساله دکتری
%        را دارید، کلمه phd را قرار دهید. کلیه تنظیمات لازم، به طور خودکار، اعمال می‌شود.
%       فایل commands.tex را حتماً به دقت مطالعه کنید؛ چون دستورات مربوط به فراخوانی بسته زی‌پرشین 
%       و دیگر بسته‌ها و ... در این فایل قرار دارد و بهتر است که با نحوه استفاده از آنها آشنا شوید.
\usepackage{amsmath, amssymb, amsthm} % بسته‌های ریاضی و قضایای امسال
\usepackage[all]{xy} % برای ترسیم نمودارها
\usepackage[left=1.3cm, right=1.3cm, top=1cm, bottom=1.3cm, nohead, hmargin=2cm,vmargin=2.5cm]{geometry} % تنظیمات حاشیه‌ها
\usepackage{graphicx} % برای درج تصاویر
\usepackage{fancyhdr} % برای تنظیم سربرگ و پاورقی صفحات
\usepackage{fontspec} % برای انتخاب فونتها
\usepackage{framed} % برای ایجاد قاب برای متن
\usepackage{ } % برای شمارش تعداد صفحات
\usepackage{cite} % برای ارجاع به منابع
\usepackage{pdfpages} % برای درج صفحات PDF
\usepackage{ifthen} % برای شرطی کردن تنظیمات
\usepackage{xcolor, colortbl} % برای استفاده از رنگها
\usepackage{pifont} % برای استفاده از نمادها

\usepackage{enumitem} % برای سفارشی‌سازی محیط‌های شماره‌گذاری
\usepackage{stackengine} % برای ترتیب متون و اشیاء
\usepackage{tabularx} % برای جداول با عرض متغیر
\usepackage{multirow} % برای ترکیب سلول‌های چند ردیفی در جدول
\usepackage{fourier} % برای فونتهای تایمز
\usepackage{array} % برای سفارشی‌سازی جداول
\usepackage{makecell} % برای سفارشی‌سازی سلول‌ها در جدول
\usepackage{ptext} % تولید متن تصادفی
\usepackage{hyperref} % برای ایجاد لینکهای رنگی
\usepackage[nottoc]{tocbibind} % برای اضافه کردن فهرست مطالب به فهرست مطالب
\usepackage{bidipoem} % برای شعر فارسی
\usepackage{tikz} % برای ترسیم نمودارهای تعاملی
\usepackage{atbegshi} % برای کنترل تصاویر و نمودارها
\usetikzlibrary{calc} % کتابخانه‌ای برای ترسیم نمودارهای TikZ


\usepackage{fancyhdr} % Custom headers and footers


% چنانچه قصد پرینت گرفتن نوشته خود را دارید، خط بالا را غیرفعال و  از دستور زیر استفاده کنید چون در صورت استفاده از دستور زیر‌‌، 
% لینک‌ها به رنگ سیاه ظاهر خواهند شد که برای پرینت گرفتن، مناسب‌تر است
%\usepackage[pagebackref=false]{hyperref}


% بسته‌ لازم برای تنظیم سربرگ‌ها
% بسته‌ای برای ظاهر شدن «مراجع» و «نمایه» در فهرست مطالب

% دستورات مربوط به ایجاد نمایه
%\usepackage{makeidx}
%\makeindex
%%%%%%%%%%%%%%%%%%%%%%%%%%
% فراخوانی بسته زی‌پرشین و تعریف قلم فارسی و انگلیسی
\usepackage{xepersian}
\settextfont[Scale=1]{XB Niloofar}
%%%%%%%%%%%%%%%%%%%%%%%%%%
% چنانچه می‌خواهید اعداد در فرمول‌ها، انگلیسی باشد، خط زیر را غیرفعال کنید
\setdigitfont[Scale=1.2]{XB Niloofar}

\defpersianfont\nas[Scale=2]{IranNastaliq}

\defpersianfont\titr[Scale=1.5]{B Titr Bold}

%%%%%%%%%%%%%%%%%%%%%%%%%%


%%%%%%%%%%%%%%%%%%%%%%%%%%
% تعریف و نحوه ظاهر شدن عنوان قضیه‌ها، تعریف‌ها، مثال‌ها و ...
\theoremstyle{definition}
\newtheorem{definition}{تعریف}[section]
\theoremstyle{theorem}
%\newtheorem{s}[definition]{}
\newtheorem{theorem}[definition]{قضیه}
\newtheorem{lemma}[definition]{لم}
\newtheorem{proposition}[definition]{گزاره}
\newtheorem{corollary}[definition]{نتیجه}
\newtheorem{remark}[definition]{تبصره}
\theoremstyle{definition}
\newtheorem{example}[definition]{مثال}
\newtheorem{notation}[definition]{نمادگذاری}
\newtheorem{definitiont}[definition]{تعریف و نمادگذاری}
\newtheorem{obs}[definition]{ملاحظه}
%%%%%%%%%%%%%%%%%%%%%%%%%%%%

\begin{document}
\baselineskip=1cm
\pagenumbering{harfi}

\newpage\clearpage
%%%%%%%%%%%%%%%%%%%%%%%%%%%%%%%%%%%%
% نیایش

\begin{center}
	\includegraphics[width=0.6\textwidth]{./BaseFile/allah.jpg}
\end{center}
\baselineskip=.750cm
\thispagestyle{empty}

% در این فایل، عنوان پایان‌نامه، مشخصات خود، متن تقدیمی‌، ستایش، سپاس‌گزاری و چکیده پایان‌نامه را به فارسی، وارد کنید.
% توجه داشته باشید که جدول حاوی مشخصات پروژه/پایان‌نامه/رساله و همچنین، مشخصات داخل آن، به طور خودکار، درج می‌شود.
%%%%%%%%%%%%%%%%%%%%%%%%%%%%%%%%%%%%


% با تعریف نام و نام خانوادگی پژوهشگر در دو متغیر زیر نام و نام خانوادگی به صورت خودکار در مکان های مورد نیاز قرارداده خواهد شد 
\newcommand{\fullname}{نام و نام خانوادگی}
\newcommand{\fullnameEN}{First and last name}
% دانشگاه خود را وارد کنید
\begin{titlepage}
\begin{center}
\vspace{-2cm}
\centerline{{\includegraphics[height=4cm]{./BaseFile/logo.jpg}}}

\vspace{-1cm}
{\nazzz \textbf{دانشگاه اصفهان}}
\vspace{-.3cm}

{\nazz \textbf{دانشکده ریاضی و آمار}}
\vspace{-.3cm}

{\naz \textbf{گروه ریاضی کاربردی و علوم کامپیوتر}}

\vspace{1.5cm}
\textbf{\nazzzzz گرایش }

\vspace{1.5cm}
\textbf{\Titr عنوان }


\vspace{2cm}
\textbf{\nazzzz{استاد راهنما:}}
\vspace{.2cm}

\textbf{\nazzz {دکتر ...}}

\textbf{\nazzz {دکتر ...}}

\vspace{1.5cm}

\textbf{\nazzzz{استاد مشاور:}}

\vspace{.2cm}

\textit{\nazzz {دکتر ...}}

\vspace{1.5cm}

\textbf{\nazzzz{پژوهشگر:}}
\vspace{.2cm}

\textbf{\nazzz {\fullname}}


\vspace{1cm}
\textbf{\nazzz{خرداد 1404}}
\end{center}
\end{titlepage}
\newpage

%%%%%%%%%%%%%%%%%%%%%%%%%%%%%%%%%%%%
\thispagestyle{empty}
% سپاس‌گزاری
\noindent{\nastaliq
	سپاس گزاری...
}\\[2cm]
{\large
	
}
\newpage
\thispagestyle{empty}
% تقدیم به
\noindent{\nastaliq
	تقدیم به ...
}
% با استفاده از دستور زیر، امضای شما، به طور خودکار، درج می‌شود.
%\signature
%\newpage\clearpage
%%%%%%%%%%%%%%%%%%%%%%%%%%%%%%%%%%%%
\newpage
\vspace*{1.7cm}
\thispagestyle{empty}
\begin{center}
	 تعهد نامه اصالت اثر
 \end{center}
 \noindent 
اینجانب
\textbf{\fullname} 
دانشجوی مقطع  دکتری 
$\square$   
کارشناسی ارشد
$\blacksquare$   
 رشته
علوم کامپیوتر
گرایش
 محاسبات نرم و هوش مصنوعی
متعهد می شوم که مطالب مندرج در این رساله/پایان‌نامه و بروندادهای منتشره در این رابطه حاصل کار پژوهشی اینجانب است و به دستاوردهای پژوهشی دیگران که در اين پژوهش از آنها استفاده شده است، مطابق مقررات ارجاع و در فهرست منابع و ماخذ ذکر گردیده است. این رساله/ پایان‌نامه قبلاً برای احراز هیچ مدرک هم سطح يا بالاتر ارائه نشده است. لذا در صورت اثبات تخلف (در هر زمان) دانشگاه اصفهان حق دارد مدرک تحصیلی صادر شده برای اینجانب را از اعتبار ساقط و ضمن درج عام موضوع در جراید کثیر الانتشار، کلیه امتیازات و حقوقی که به موجب آن در طی دوره تحصیل و مدت زمان بعد از فراغت از تحصیل تا اثبات تخلف به ذینفعان تعلق گرفته را مسترد گرداند.\\ 
کلیه حقوق مادی و معنوی این اثر متعلق به دانشگاه اصفهان می باشد.
\begin{table}[h!]
	\renewcommand\arraystretch{3} 	
	\begin{tabular}{rr}
	نام و نام خانوادگی دانشجو: 
		& \hspace{5cm}  امضاء
		   \\
	تأیید استاد (استادان) راهنما: نام و نام خانوادگی 
		&  \hspace{5cm} امضاء
		  \\
	\end{tabular}
\end{table}

%%%%%%%%%%%%%%%%%%%%%%%%%%%%%%%%%%%%.

 
\chapter*{چکیده}
\ptext[1-2]


\tableofcontents
\ptext[3-6]
\listoffigures
\newpage
\pagenumbering{arabic}
\chapter{مقدمه}
\documentclass[a4paper 11px twocolumn]{article}
% تنظیمات و پکیج های مورد نیاز در این فایل قرار دارد. برای تنظیم و افزودن پکیج به این فایل مراجعه کنید.
\usepackage{amsmath, amssymb, amsthm} % بسته‌های ریاضی و قضایای امسال
\usepackage[all]{xy} % برای ترسیم نمودارها
\usepackage[left=1.3cm, right=1.3cm, top=1cm, bottom=1.3cm, nohead, hmargin=2cm,vmargin=2.5cm]{geometry} % تنظیمات حاشیه‌ها
\usepackage{graphicx} % برای درج تصاویر
\usepackage{fancyhdr} % برای تنظیم سربرگ و پاورقی صفحات
\usepackage{fontspec} % برای انتخاب فونتها
\usepackage{framed} % برای ایجاد قاب برای متن
\usepackage{ } % برای شمارش تعداد صفحات
\usepackage{cite} % برای ارجاع به منابع
\usepackage{pdfpages} % برای درج صفحات PDF
\usepackage{ifthen} % برای شرطی کردن تنظیمات
\usepackage{xcolor, colortbl} % برای استفاده از رنگها
\usepackage{pifont} % برای استفاده از نمادها

\usepackage{enumitem} % برای سفارشی‌سازی محیط‌های شماره‌گذاری
\usepackage{stackengine} % برای ترتیب متون و اشیاء
\usepackage{tabularx} % برای جداول با عرض متغیر
\usepackage{multirow} % برای ترکیب سلول‌های چند ردیفی در جدول
\usepackage{fourier} % برای فونتهای تایمز
\usepackage{array} % برای سفارشی‌سازی جداول
\usepackage{makecell} % برای سفارشی‌سازی سلول‌ها در جدول
\usepackage{ptext} % تولید متن تصادفی
\usepackage{hyperref} % برای ایجاد لینکهای رنگی
\usepackage[nottoc]{tocbibind} % برای اضافه کردن فهرست مطالب به فهرست مطالب
\usepackage{bidipoem} % برای شعر فارسی
\usepackage{tikz} % برای ترسیم نمودارهای تعاملی
\usepackage{atbegshi} % برای کنترل تصاویر و نمودارها
\usetikzlibrary{calc} % کتابخانه‌ای برای ترسیم نمودارهای TikZ


\usepackage{fancyhdr} % Custom headers and footers


% چنانچه قصد پرینت گرفتن نوشته خود را دارید، خط بالا را غیرفعال و  از دستور زیر استفاده کنید چون در صورت استفاده از دستور زیر‌‌، 
% لینک‌ها به رنگ سیاه ظاهر خواهند شد که برای پرینت گرفتن، مناسب‌تر است
%\usepackage[pagebackref=false]{hyperref}


% بسته‌ لازم برای تنظیم سربرگ‌ها
% بسته‌ای برای ظاهر شدن «مراجع» و «نمایه» در فهرست مطالب

% دستورات مربوط به ایجاد نمایه
%\usepackage{makeidx}
%\makeindex
%%%%%%%%%%%%%%%%%%%%%%%%%%
% فراخوانی بسته زی‌پرشین و تعریف قلم فارسی و انگلیسی
\usepackage{xepersian}
\settextfont[Scale=1]{XB Niloofar}
%%%%%%%%%%%%%%%%%%%%%%%%%%
% چنانچه می‌خواهید اعداد در فرمول‌ها، انگلیسی باشد، خط زیر را غیرفعال کنید
\setdigitfont[Scale=1.2]{XB Niloofar}

\defpersianfont\nas[Scale=2]{IranNastaliq}

\defpersianfont\titr[Scale=1.5]{B Titr Bold}

%%%%%%%%%%%%%%%%%%%%%%%%%%


%%%%%%%%%%%%%%%%%%%%%%%%%%
% تعریف و نحوه ظاهر شدن عنوان قضیه‌ها، تعریف‌ها، مثال‌ها و ...
\theoremstyle{definition}
\newtheorem{definition}{تعریف}[section]
\theoremstyle{theorem}
%\newtheorem{s}[definition]{}
\newtheorem{theorem}[definition]{قضیه}
\newtheorem{lemma}[definition]{لم}
\newtheorem{proposition}[definition]{گزاره}
\newtheorem{corollary}[definition]{نتیجه}
\newtheorem{remark}[definition]{تبصره}
\theoremstyle{definition}
\newtheorem{example}[definition]{مثال}
\newtheorem{notation}[definition]{نمادگذاری}
\newtheorem{definitiont}[definition]{تعریف و نمادگذاری}
\newtheorem{obs}[definition]{ملاحظه}
%%%%%%%%%%%%%%%%%%%%%%%%%%%%

\begin{document}
% مشخصات گزارش را در این قسمت تغییر بدهید.
\begin{titlepage}
	\setlength{\arrayrulewidth}{0.4mm}
	\newcommand{\HRule}{\rule{\linewidth}{0.5mm}}
	\begin{center}
		% logo
		\includegraphics[width=0.5\textwidth]{./Fig/logo.png}~\\[4cm]
		
		\HRule \\[0.4cm]
		{
			\Huge\textbf{عنوان گزارش} \\
			[0.4cm]
			
			\large \bfseries
			نام درس \\ 
			[0.4cm]}
		\HRule \\[4cm]
		
		\large\textbf{نویسنده: }\\[0.2cm]
		نام و نام خوانوادگی 
		\\[3cm]
		
		\vfill
		
		{\large \today}
		
	\end{center}
\end{titlepage}

% متن مربوط به گزارش را در این قسمت قرار دهید.
\twocolumn
\ptext[0-10]


% در صورت نیاز به مراجع و پیوست از دو خط و فایل های مربوطه استفاده فرمایید.
% پیوست خود را در این قسمت وارد کنید.
\appendix
\onecolumn
% شروع محیط پیوست 
\section{پیوست}
\ptext[10-20]

% مراجع خود را در این قسمت وارد کنید.
% دستوری برای کوچک کردن اندازه فونت‌ها 
\small
% شروع محیط مراجع 

\renewcommand\bibname{مراجع}
\setLTRbibitems
% برای ارجاع دهی خودکار محتوای فایل bibtex مراجع را در فایل Refrence/myreferences.bib قرار داده و ذخیره کنید.  و یا میتوانید دو خط زیر را کامنت کرده و از دستورات بعدی ان بزرای ارجاع دهی استفاده کنید.
\bibliographystyle{IEEEtran}
\bibliography{Refrence/myreferences.bib}
%\begin{thebibliography}{99}
%
%\resetlatinfont
%% تذکر: چنانچه از نسخه‌های قدیمی زی‌پرشین و بی‌دی استفاده می‌کنید، هنگام استفاده از سه دستور بالا، با خطا 
%% مواجه می‌شوید. بنابراین باید زی‌پرشین و بی‌دی خود را آپدیت کنید.
%% چنانچه مرجع فارسی هم دارید باید یا از بسته Persian-bib استفاده کنید و یا راهنمای bidi را مطالعه کنید.
%
%\bibitem {TE} { \sc S. T. Aldrich, E.  Enochs, J. R. Garc\'{i}a Rozas, L. Oyonarte,} {\sl Covers and envelopes in Grothendieck categories: flat covers of complexes with applications,} J. Algebra {\bf 243} (2001), 615-630.
%
%\bibitem {AJL}{ \sc L. Alonso, A. Jerem\'{i}oas, J. Lipman,} {\sl  Duality and flat base change on formal schemes,}
%in “Studies in Duality on Noetherian Formal Schemes and Non-Noetherian Ordinary
%Schemes,” Contemp. Math. {\bf 244}, Amer. Math. Soc., 1999, 3-90.
%
%\bibitem {AF} {\sc F. W. Anderson, K. R. Fuller,} {\sl Rings and categories of modules,} Graduate Texts in Mathematics 13, Springer-Verlag, Berlin-Heidelberg-New York, 1974.
%
%\bibitem {AEHS} {\sc J. Asadollahi, H. Eshraghi, R. Hafezi, Sh. Salarian,} {\sl On the homotopy categories of projective and injective representations of quivers,} J. Algebra {\bf 346} (2011), 101-115.
%
%
%
%
%\end{thebibliography}
%
%

\end{document}
\chapter{مرور ادبیات و پیشینه تحقیق}
\documentclass[a4paper 11px twocolumn]{article}
% تنظیمات و پکیج های مورد نیاز در این فایل قرار دارد. برای تنظیم و افزودن پکیج به این فایل مراجعه کنید.
\usepackage{amsmath, amssymb, amsthm} % بسته‌های ریاضی و قضایای امسال
\usepackage[all]{xy} % برای ترسیم نمودارها
\usepackage[left=1.3cm, right=1.3cm, top=1cm, bottom=1.3cm, nohead, hmargin=2cm,vmargin=2.5cm]{geometry} % تنظیمات حاشیه‌ها
\usepackage{graphicx} % برای درج تصاویر
\usepackage{fancyhdr} % برای تنظیم سربرگ و پاورقی صفحات
\usepackage{fontspec} % برای انتخاب فونتها
\usepackage{framed} % برای ایجاد قاب برای متن
\usepackage{ } % برای شمارش تعداد صفحات
\usepackage{cite} % برای ارجاع به منابع
\usepackage{pdfpages} % برای درج صفحات PDF
\usepackage{ifthen} % برای شرطی کردن تنظیمات
\usepackage{xcolor, colortbl} % برای استفاده از رنگها
\usepackage{pifont} % برای استفاده از نمادها

\usepackage{enumitem} % برای سفارشی‌سازی محیط‌های شماره‌گذاری
\usepackage{stackengine} % برای ترتیب متون و اشیاء
\usepackage{tabularx} % برای جداول با عرض متغیر
\usepackage{multirow} % برای ترکیب سلول‌های چند ردیفی در جدول
\usepackage{fourier} % برای فونتهای تایمز
\usepackage{array} % برای سفارشی‌سازی جداول
\usepackage{makecell} % برای سفارشی‌سازی سلول‌ها در جدول
\usepackage{ptext} % تولید متن تصادفی
\usepackage{hyperref} % برای ایجاد لینکهای رنگی
\usepackage[nottoc]{tocbibind} % برای اضافه کردن فهرست مطالب به فهرست مطالب
\usepackage{bidipoem} % برای شعر فارسی
\usepackage{tikz} % برای ترسیم نمودارهای تعاملی
\usepackage{atbegshi} % برای کنترل تصاویر و نمودارها
\usetikzlibrary{calc} % کتابخانه‌ای برای ترسیم نمودارهای TikZ


\usepackage{fancyhdr} % Custom headers and footers


% چنانچه قصد پرینت گرفتن نوشته خود را دارید، خط بالا را غیرفعال و  از دستور زیر استفاده کنید چون در صورت استفاده از دستور زیر‌‌، 
% لینک‌ها به رنگ سیاه ظاهر خواهند شد که برای پرینت گرفتن، مناسب‌تر است
%\usepackage[pagebackref=false]{hyperref}


% بسته‌ لازم برای تنظیم سربرگ‌ها
% بسته‌ای برای ظاهر شدن «مراجع» و «نمایه» در فهرست مطالب

% دستورات مربوط به ایجاد نمایه
%\usepackage{makeidx}
%\makeindex
%%%%%%%%%%%%%%%%%%%%%%%%%%
% فراخوانی بسته زی‌پرشین و تعریف قلم فارسی و انگلیسی
\usepackage{xepersian}
\settextfont[Scale=1]{XB Niloofar}
%%%%%%%%%%%%%%%%%%%%%%%%%%
% چنانچه می‌خواهید اعداد در فرمول‌ها، انگلیسی باشد، خط زیر را غیرفعال کنید
\setdigitfont[Scale=1.2]{XB Niloofar}

\defpersianfont\nas[Scale=2]{IranNastaliq}

\defpersianfont\titr[Scale=1.5]{B Titr Bold}

%%%%%%%%%%%%%%%%%%%%%%%%%%


%%%%%%%%%%%%%%%%%%%%%%%%%%
% تعریف و نحوه ظاهر شدن عنوان قضیه‌ها، تعریف‌ها، مثال‌ها و ...
\theoremstyle{definition}
\newtheorem{definition}{تعریف}[section]
\theoremstyle{theorem}
%\newtheorem{s}[definition]{}
\newtheorem{theorem}[definition]{قضیه}
\newtheorem{lemma}[definition]{لم}
\newtheorem{proposition}[definition]{گزاره}
\newtheorem{corollary}[definition]{نتیجه}
\newtheorem{remark}[definition]{تبصره}
\theoremstyle{definition}
\newtheorem{example}[definition]{مثال}
\newtheorem{notation}[definition]{نمادگذاری}
\newtheorem{definitiont}[definition]{تعریف و نمادگذاری}
\newtheorem{obs}[definition]{ملاحظه}
%%%%%%%%%%%%%%%%%%%%%%%%%%%%

\begin{document}
% مشخصات گزارش را در این قسمت تغییر بدهید.
\begin{titlepage}
	\setlength{\arrayrulewidth}{0.4mm}
	\newcommand{\HRule}{\rule{\linewidth}{0.5mm}}
	\begin{center}
		% logo
		\includegraphics[width=0.5\textwidth]{./Fig/logo.png}~\\[4cm]
		
		\HRule \\[0.4cm]
		{
			\Huge\textbf{عنوان گزارش} \\
			[0.4cm]
			
			\large \bfseries
			نام درس \\ 
			[0.4cm]}
		\HRule \\[4cm]
		
		\large\textbf{نویسنده: }\\[0.2cm]
		نام و نام خوانوادگی 
		\\[3cm]
		
		\vfill
		
		{\large \today}
		
	\end{center}
\end{titlepage}

% متن مربوط به گزارش را در این قسمت قرار دهید.
\twocolumn
\ptext[0-10]


% در صورت نیاز به مراجع و پیوست از دو خط و فایل های مربوطه استفاده فرمایید.
% پیوست خود را در این قسمت وارد کنید.
\appendix
\onecolumn
% شروع محیط پیوست 
\section{پیوست}
\ptext[10-20]

% مراجع خود را در این قسمت وارد کنید.
% دستوری برای کوچک کردن اندازه فونت‌ها 
\small
% شروع محیط مراجع 

\renewcommand\bibname{مراجع}
\setLTRbibitems
% برای ارجاع دهی خودکار محتوای فایل bibtex مراجع را در فایل Refrence/myreferences.bib قرار داده و ذخیره کنید.  و یا میتوانید دو خط زیر را کامنت کرده و از دستورات بعدی ان بزرای ارجاع دهی استفاده کنید.
\bibliographystyle{IEEEtran}
\bibliography{Refrence/myreferences.bib}
%\begin{thebibliography}{99}
%
%\resetlatinfont
%% تذکر: چنانچه از نسخه‌های قدیمی زی‌پرشین و بی‌دی استفاده می‌کنید، هنگام استفاده از سه دستور بالا، با خطا 
%% مواجه می‌شوید. بنابراین باید زی‌پرشین و بی‌دی خود را آپدیت کنید.
%% چنانچه مرجع فارسی هم دارید باید یا از بسته Persian-bib استفاده کنید و یا راهنمای bidi را مطالعه کنید.
%
%\bibitem {TE} { \sc S. T. Aldrich, E.  Enochs, J. R. Garc\'{i}a Rozas, L. Oyonarte,} {\sl Covers and envelopes in Grothendieck categories: flat covers of complexes with applications,} J. Algebra {\bf 243} (2001), 615-630.
%
%\bibitem {AJL}{ \sc L. Alonso, A. Jerem\'{i}oas, J. Lipman,} {\sl  Duality and flat base change on formal schemes,}
%in “Studies in Duality on Noetherian Formal Schemes and Non-Noetherian Ordinary
%Schemes,” Contemp. Math. {\bf 244}, Amer. Math. Soc., 1999, 3-90.
%
%\bibitem {AF} {\sc F. W. Anderson, K. R. Fuller,} {\sl Rings and categories of modules,} Graduate Texts in Mathematics 13, Springer-Verlag, Berlin-Heidelberg-New York, 1974.
%
%\bibitem {AEHS} {\sc J. Asadollahi, H. Eshraghi, R. Hafezi, Sh. Salarian,} {\sl On the homotopy categories of projective and injective representations of quivers,} J. Algebra {\bf 346} (2011), 101-115.
%
%
%
%
%\end{thebibliography}
%
%

\end{document}
\chapter{مدل‌ها، چالش‌ها و ارزیابی}
\documentclass[a4paper 11px twocolumn]{article}
% تنظیمات و پکیج های مورد نیاز در این فایل قرار دارد. برای تنظیم و افزودن پکیج به این فایل مراجعه کنید.
\usepackage{amsmath, amssymb, amsthm} % بسته‌های ریاضی و قضایای امسال
\usepackage[all]{xy} % برای ترسیم نمودارها
\usepackage[left=1.3cm, right=1.3cm, top=1cm, bottom=1.3cm, nohead, hmargin=2cm,vmargin=2.5cm]{geometry} % تنظیمات حاشیه‌ها
\usepackage{graphicx} % برای درج تصاویر
\usepackage{fancyhdr} % برای تنظیم سربرگ و پاورقی صفحات
\usepackage{fontspec} % برای انتخاب فونتها
\usepackage{framed} % برای ایجاد قاب برای متن
\usepackage{ } % برای شمارش تعداد صفحات
\usepackage{cite} % برای ارجاع به منابع
\usepackage{pdfpages} % برای درج صفحات PDF
\usepackage{ifthen} % برای شرطی کردن تنظیمات
\usepackage{xcolor, colortbl} % برای استفاده از رنگها
\usepackage{pifont} % برای استفاده از نمادها

\usepackage{enumitem} % برای سفارشی‌سازی محیط‌های شماره‌گذاری
\usepackage{stackengine} % برای ترتیب متون و اشیاء
\usepackage{tabularx} % برای جداول با عرض متغیر
\usepackage{multirow} % برای ترکیب سلول‌های چند ردیفی در جدول
\usepackage{fourier} % برای فونتهای تایمز
\usepackage{array} % برای سفارشی‌سازی جداول
\usepackage{makecell} % برای سفارشی‌سازی سلول‌ها در جدول
\usepackage{ptext} % تولید متن تصادفی
\usepackage{hyperref} % برای ایجاد لینکهای رنگی
\usepackage[nottoc]{tocbibind} % برای اضافه کردن فهرست مطالب به فهرست مطالب
\usepackage{bidipoem} % برای شعر فارسی
\usepackage{tikz} % برای ترسیم نمودارهای تعاملی
\usepackage{atbegshi} % برای کنترل تصاویر و نمودارها
\usetikzlibrary{calc} % کتابخانه‌ای برای ترسیم نمودارهای TikZ


\usepackage{fancyhdr} % Custom headers and footers


% چنانچه قصد پرینت گرفتن نوشته خود را دارید، خط بالا را غیرفعال و  از دستور زیر استفاده کنید چون در صورت استفاده از دستور زیر‌‌، 
% لینک‌ها به رنگ سیاه ظاهر خواهند شد که برای پرینت گرفتن، مناسب‌تر است
%\usepackage[pagebackref=false]{hyperref}


% بسته‌ لازم برای تنظیم سربرگ‌ها
% بسته‌ای برای ظاهر شدن «مراجع» و «نمایه» در فهرست مطالب

% دستورات مربوط به ایجاد نمایه
%\usepackage{makeidx}
%\makeindex
%%%%%%%%%%%%%%%%%%%%%%%%%%
% فراخوانی بسته زی‌پرشین و تعریف قلم فارسی و انگلیسی
\usepackage{xepersian}
\settextfont[Scale=1]{XB Niloofar}
%%%%%%%%%%%%%%%%%%%%%%%%%%
% چنانچه می‌خواهید اعداد در فرمول‌ها، انگلیسی باشد، خط زیر را غیرفعال کنید
\setdigitfont[Scale=1.2]{XB Niloofar}

\defpersianfont\nas[Scale=2]{IranNastaliq}

\defpersianfont\titr[Scale=1.5]{B Titr Bold}

%%%%%%%%%%%%%%%%%%%%%%%%%%


%%%%%%%%%%%%%%%%%%%%%%%%%%
% تعریف و نحوه ظاهر شدن عنوان قضیه‌ها، تعریف‌ها، مثال‌ها و ...
\theoremstyle{definition}
\newtheorem{definition}{تعریف}[section]
\theoremstyle{theorem}
%\newtheorem{s}[definition]{}
\newtheorem{theorem}[definition]{قضیه}
\newtheorem{lemma}[definition]{لم}
\newtheorem{proposition}[definition]{گزاره}
\newtheorem{corollary}[definition]{نتیجه}
\newtheorem{remark}[definition]{تبصره}
\theoremstyle{definition}
\newtheorem{example}[definition]{مثال}
\newtheorem{notation}[definition]{نمادگذاری}
\newtheorem{definitiont}[definition]{تعریف و نمادگذاری}
\newtheorem{obs}[definition]{ملاحظه}
%%%%%%%%%%%%%%%%%%%%%%%%%%%%

\begin{document}
% مشخصات گزارش را در این قسمت تغییر بدهید.
\begin{titlepage}
	\setlength{\arrayrulewidth}{0.4mm}
	\newcommand{\HRule}{\rule{\linewidth}{0.5mm}}
	\begin{center}
		% logo
		\includegraphics[width=0.5\textwidth]{./Fig/logo.png}~\\[4cm]
		
		\HRule \\[0.4cm]
		{
			\Huge\textbf{عنوان گزارش} \\
			[0.4cm]
			
			\large \bfseries
			نام درس \\ 
			[0.4cm]}
		\HRule \\[4cm]
		
		\large\textbf{نویسنده: }\\[0.2cm]
		نام و نام خوانوادگی 
		\\[3cm]
		
		\vfill
		
		{\large \today}
		
	\end{center}
\end{titlepage}

% متن مربوط به گزارش را در این قسمت قرار دهید.
\twocolumn
\ptext[0-10]


% در صورت نیاز به مراجع و پیوست از دو خط و فایل های مربوطه استفاده فرمایید.
% پیوست خود را در این قسمت وارد کنید.
\appendix
\onecolumn
% شروع محیط پیوست 
\section{پیوست}
\ptext[10-20]

% مراجع خود را در این قسمت وارد کنید.
% دستوری برای کوچک کردن اندازه فونت‌ها 
\small
% شروع محیط مراجع 

\renewcommand\bibname{مراجع}
\setLTRbibitems
% برای ارجاع دهی خودکار محتوای فایل bibtex مراجع را در فایل Refrence/myreferences.bib قرار داده و ذخیره کنید.  و یا میتوانید دو خط زیر را کامنت کرده و از دستورات بعدی ان بزرای ارجاع دهی استفاده کنید.
\bibliographystyle{IEEEtran}
\bibliography{Refrence/myreferences.bib}
%\begin{thebibliography}{99}
%
%\resetlatinfont
%% تذکر: چنانچه از نسخه‌های قدیمی زی‌پرشین و بی‌دی استفاده می‌کنید، هنگام استفاده از سه دستور بالا، با خطا 
%% مواجه می‌شوید. بنابراین باید زی‌پرشین و بی‌دی خود را آپدیت کنید.
%% چنانچه مرجع فارسی هم دارید باید یا از بسته Persian-bib استفاده کنید و یا راهنمای bidi را مطالعه کنید.
%
%\bibitem {TE} { \sc S. T. Aldrich, E.  Enochs, J. R. Garc\'{i}a Rozas, L. Oyonarte,} {\sl Covers and envelopes in Grothendieck categories: flat covers of complexes with applications,} J. Algebra {\bf 243} (2001), 615-630.
%
%\bibitem {AJL}{ \sc L. Alonso, A. Jerem\'{i}oas, J. Lipman,} {\sl  Duality and flat base change on formal schemes,}
%in “Studies in Duality on Noetherian Formal Schemes and Non-Noetherian Ordinary
%Schemes,” Contemp. Math. {\bf 244}, Amer. Math. Soc., 1999, 3-90.
%
%\bibitem {AF} {\sc F. W. Anderson, K. R. Fuller,} {\sl Rings and categories of modules,} Graduate Texts in Mathematics 13, Springer-Verlag, Berlin-Heidelberg-New York, 1974.
%
%\bibitem {AEHS} {\sc J. Asadollahi, H. Eshraghi, R. Hafezi, Sh. Salarian,} {\sl On the homotopy categories of projective and injective representations of quivers,} J. Algebra {\bf 346} (2011), 101-115.
%
%
%
%
%\end{thebibliography}
%
%

\end{document}
\chapter{نتایج و بحث}
\documentclass[a4paper 11px twocolumn]{article}
% تنظیمات و پکیج های مورد نیاز در این فایل قرار دارد. برای تنظیم و افزودن پکیج به این فایل مراجعه کنید.
\usepackage{amsmath, amssymb, amsthm} % بسته‌های ریاضی و قضایای امسال
\usepackage[all]{xy} % برای ترسیم نمودارها
\usepackage[left=1.3cm, right=1.3cm, top=1cm, bottom=1.3cm, nohead, hmargin=2cm,vmargin=2.5cm]{geometry} % تنظیمات حاشیه‌ها
\usepackage{graphicx} % برای درج تصاویر
\usepackage{fancyhdr} % برای تنظیم سربرگ و پاورقی صفحات
\usepackage{fontspec} % برای انتخاب فونتها
\usepackage{framed} % برای ایجاد قاب برای متن
\usepackage{ } % برای شمارش تعداد صفحات
\usepackage{cite} % برای ارجاع به منابع
\usepackage{pdfpages} % برای درج صفحات PDF
\usepackage{ifthen} % برای شرطی کردن تنظیمات
\usepackage{xcolor, colortbl} % برای استفاده از رنگها
\usepackage{pifont} % برای استفاده از نمادها

\usepackage{enumitem} % برای سفارشی‌سازی محیط‌های شماره‌گذاری
\usepackage{stackengine} % برای ترتیب متون و اشیاء
\usepackage{tabularx} % برای جداول با عرض متغیر
\usepackage{multirow} % برای ترکیب سلول‌های چند ردیفی در جدول
\usepackage{fourier} % برای فونتهای تایمز
\usepackage{array} % برای سفارشی‌سازی جداول
\usepackage{makecell} % برای سفارشی‌سازی سلول‌ها در جدول
\usepackage{ptext} % تولید متن تصادفی
\usepackage{hyperref} % برای ایجاد لینکهای رنگی
\usepackage[nottoc]{tocbibind} % برای اضافه کردن فهرست مطالب به فهرست مطالب
\usepackage{bidipoem} % برای شعر فارسی
\usepackage{tikz} % برای ترسیم نمودارهای تعاملی
\usepackage{atbegshi} % برای کنترل تصاویر و نمودارها
\usetikzlibrary{calc} % کتابخانه‌ای برای ترسیم نمودارهای TikZ


\usepackage{fancyhdr} % Custom headers and footers


% چنانچه قصد پرینت گرفتن نوشته خود را دارید، خط بالا را غیرفعال و  از دستور زیر استفاده کنید چون در صورت استفاده از دستور زیر‌‌، 
% لینک‌ها به رنگ سیاه ظاهر خواهند شد که برای پرینت گرفتن، مناسب‌تر است
%\usepackage[pagebackref=false]{hyperref}


% بسته‌ لازم برای تنظیم سربرگ‌ها
% بسته‌ای برای ظاهر شدن «مراجع» و «نمایه» در فهرست مطالب

% دستورات مربوط به ایجاد نمایه
%\usepackage{makeidx}
%\makeindex
%%%%%%%%%%%%%%%%%%%%%%%%%%
% فراخوانی بسته زی‌پرشین و تعریف قلم فارسی و انگلیسی
\usepackage{xepersian}
\settextfont[Scale=1]{XB Niloofar}
%%%%%%%%%%%%%%%%%%%%%%%%%%
% چنانچه می‌خواهید اعداد در فرمول‌ها، انگلیسی باشد، خط زیر را غیرفعال کنید
\setdigitfont[Scale=1.2]{XB Niloofar}

\defpersianfont\nas[Scale=2]{IranNastaliq}

\defpersianfont\titr[Scale=1.5]{B Titr Bold}

%%%%%%%%%%%%%%%%%%%%%%%%%%


%%%%%%%%%%%%%%%%%%%%%%%%%%
% تعریف و نحوه ظاهر شدن عنوان قضیه‌ها، تعریف‌ها، مثال‌ها و ...
\theoremstyle{definition}
\newtheorem{definition}{تعریف}[section]
\theoremstyle{theorem}
%\newtheorem{s}[definition]{}
\newtheorem{theorem}[definition]{قضیه}
\newtheorem{lemma}[definition]{لم}
\newtheorem{proposition}[definition]{گزاره}
\newtheorem{corollary}[definition]{نتیجه}
\newtheorem{remark}[definition]{تبصره}
\theoremstyle{definition}
\newtheorem{example}[definition]{مثال}
\newtheorem{notation}[definition]{نمادگذاری}
\newtheorem{definitiont}[definition]{تعریف و نمادگذاری}
\newtheorem{obs}[definition]{ملاحظه}
%%%%%%%%%%%%%%%%%%%%%%%%%%%%

\begin{document}
% مشخصات گزارش را در این قسمت تغییر بدهید.
\begin{titlepage}
	\setlength{\arrayrulewidth}{0.4mm}
	\newcommand{\HRule}{\rule{\linewidth}{0.5mm}}
	\begin{center}
		% logo
		\includegraphics[width=0.5\textwidth]{./Fig/logo.png}~\\[4cm]
		
		\HRule \\[0.4cm]
		{
			\Huge\textbf{عنوان گزارش} \\
			[0.4cm]
			
			\large \bfseries
			نام درس \\ 
			[0.4cm]}
		\HRule \\[4cm]
		
		\large\textbf{نویسنده: }\\[0.2cm]
		نام و نام خوانوادگی 
		\\[3cm]
		
		\vfill
		
		{\large \today}
		
	\end{center}
\end{titlepage}

% متن مربوط به گزارش را در این قسمت قرار دهید.
\twocolumn
\ptext[0-10]


% در صورت نیاز به مراجع و پیوست از دو خط و فایل های مربوطه استفاده فرمایید.
% پیوست خود را در این قسمت وارد کنید.
\appendix
\onecolumn
% شروع محیط پیوست 
\section{پیوست}
\ptext[10-20]

% مراجع خود را در این قسمت وارد کنید.
% دستوری برای کوچک کردن اندازه فونت‌ها 
\small
% شروع محیط مراجع 

\renewcommand\bibname{مراجع}
\setLTRbibitems
% برای ارجاع دهی خودکار محتوای فایل bibtex مراجع را در فایل Refrence/myreferences.bib قرار داده و ذخیره کنید.  و یا میتوانید دو خط زیر را کامنت کرده و از دستورات بعدی ان بزرای ارجاع دهی استفاده کنید.
\bibliographystyle{IEEEtran}
\bibliography{Refrence/myreferences.bib}
%\begin{thebibliography}{99}
%
%\resetlatinfont
%% تذکر: چنانچه از نسخه‌های قدیمی زی‌پرشین و بی‌دی استفاده می‌کنید، هنگام استفاده از سه دستور بالا، با خطا 
%% مواجه می‌شوید. بنابراین باید زی‌پرشین و بی‌دی خود را آپدیت کنید.
%% چنانچه مرجع فارسی هم دارید باید یا از بسته Persian-bib استفاده کنید و یا راهنمای bidi را مطالعه کنید.
%
%\bibitem {TE} { \sc S. T. Aldrich, E.  Enochs, J. R. Garc\'{i}a Rozas, L. Oyonarte,} {\sl Covers and envelopes in Grothendieck categories: flat covers of complexes with applications,} J. Algebra {\bf 243} (2001), 615-630.
%
%\bibitem {AJL}{ \sc L. Alonso, A. Jerem\'{i}oas, J. Lipman,} {\sl  Duality and flat base change on formal schemes,}
%in “Studies in Duality on Noetherian Formal Schemes and Non-Noetherian Ordinary
%Schemes,” Contemp. Math. {\bf 244}, Amer. Math. Soc., 1999, 3-90.
%
%\bibitem {AF} {\sc F. W. Anderson, K. R. Fuller,} {\sl Rings and categories of modules,} Graduate Texts in Mathematics 13, Springer-Verlag, Berlin-Heidelberg-New York, 1974.
%
%\bibitem {AEHS} {\sc J. Asadollahi, H. Eshraghi, R. Hafezi, Sh. Salarian,} {\sl On the homotopy categories of projective and injective representations of quivers,} J. Algebra {\bf 346} (2011), 101-115.
%
%
%
%
%\end{thebibliography}
%
%

\end{document}
\chapter{نتیجه‌گیری و پیشنهادات}
\documentclass[a4paper 11px twocolumn]{article}
% تنظیمات و پکیج های مورد نیاز در این فایل قرار دارد. برای تنظیم و افزودن پکیج به این فایل مراجعه کنید.
\usepackage{amsmath, amssymb, amsthm} % بسته‌های ریاضی و قضایای امسال
\usepackage[all]{xy} % برای ترسیم نمودارها
\usepackage[left=1.3cm, right=1.3cm, top=1cm, bottom=1.3cm, nohead, hmargin=2cm,vmargin=2.5cm]{geometry} % تنظیمات حاشیه‌ها
\usepackage{graphicx} % برای درج تصاویر
\usepackage{fancyhdr} % برای تنظیم سربرگ و پاورقی صفحات
\usepackage{fontspec} % برای انتخاب فونتها
\usepackage{framed} % برای ایجاد قاب برای متن
\usepackage{ } % برای شمارش تعداد صفحات
\usepackage{cite} % برای ارجاع به منابع
\usepackage{pdfpages} % برای درج صفحات PDF
\usepackage{ifthen} % برای شرطی کردن تنظیمات
\usepackage{xcolor, colortbl} % برای استفاده از رنگها
\usepackage{pifont} % برای استفاده از نمادها

\usepackage{enumitem} % برای سفارشی‌سازی محیط‌های شماره‌گذاری
\usepackage{stackengine} % برای ترتیب متون و اشیاء
\usepackage{tabularx} % برای جداول با عرض متغیر
\usepackage{multirow} % برای ترکیب سلول‌های چند ردیفی در جدول
\usepackage{fourier} % برای فونتهای تایمز
\usepackage{array} % برای سفارشی‌سازی جداول
\usepackage{makecell} % برای سفارشی‌سازی سلول‌ها در جدول
\usepackage{ptext} % تولید متن تصادفی
\usepackage{hyperref} % برای ایجاد لینکهای رنگی
\usepackage[nottoc]{tocbibind} % برای اضافه کردن فهرست مطالب به فهرست مطالب
\usepackage{bidipoem} % برای شعر فارسی
\usepackage{tikz} % برای ترسیم نمودارهای تعاملی
\usepackage{atbegshi} % برای کنترل تصاویر و نمودارها
\usetikzlibrary{calc} % کتابخانه‌ای برای ترسیم نمودارهای TikZ


\usepackage{fancyhdr} % Custom headers and footers


% چنانچه قصد پرینت گرفتن نوشته خود را دارید، خط بالا را غیرفعال و  از دستور زیر استفاده کنید چون در صورت استفاده از دستور زیر‌‌، 
% لینک‌ها به رنگ سیاه ظاهر خواهند شد که برای پرینت گرفتن، مناسب‌تر است
%\usepackage[pagebackref=false]{hyperref}


% بسته‌ لازم برای تنظیم سربرگ‌ها
% بسته‌ای برای ظاهر شدن «مراجع» و «نمایه» در فهرست مطالب

% دستورات مربوط به ایجاد نمایه
%\usepackage{makeidx}
%\makeindex
%%%%%%%%%%%%%%%%%%%%%%%%%%
% فراخوانی بسته زی‌پرشین و تعریف قلم فارسی و انگلیسی
\usepackage{xepersian}
\settextfont[Scale=1]{XB Niloofar}
%%%%%%%%%%%%%%%%%%%%%%%%%%
% چنانچه می‌خواهید اعداد در فرمول‌ها، انگلیسی باشد، خط زیر را غیرفعال کنید
\setdigitfont[Scale=1.2]{XB Niloofar}

\defpersianfont\nas[Scale=2]{IranNastaliq}

\defpersianfont\titr[Scale=1.5]{B Titr Bold}

%%%%%%%%%%%%%%%%%%%%%%%%%%


%%%%%%%%%%%%%%%%%%%%%%%%%%
% تعریف و نحوه ظاهر شدن عنوان قضیه‌ها، تعریف‌ها، مثال‌ها و ...
\theoremstyle{definition}
\newtheorem{definition}{تعریف}[section]
\theoremstyle{theorem}
%\newtheorem{s}[definition]{}
\newtheorem{theorem}[definition]{قضیه}
\newtheorem{lemma}[definition]{لم}
\newtheorem{proposition}[definition]{گزاره}
\newtheorem{corollary}[definition]{نتیجه}
\newtheorem{remark}[definition]{تبصره}
\theoremstyle{definition}
\newtheorem{example}[definition]{مثال}
\newtheorem{notation}[definition]{نمادگذاری}
\newtheorem{definitiont}[definition]{تعریف و نمادگذاری}
\newtheorem{obs}[definition]{ملاحظه}
%%%%%%%%%%%%%%%%%%%%%%%%%%%%

\begin{document}
% مشخصات گزارش را در این قسمت تغییر بدهید.
\begin{titlepage}
	\setlength{\arrayrulewidth}{0.4mm}
	\newcommand{\HRule}{\rule{\linewidth}{0.5mm}}
	\begin{center}
		% logo
		\includegraphics[width=0.5\textwidth]{./Fig/logo.png}~\\[4cm]
		
		\HRule \\[0.4cm]
		{
			\Huge\textbf{عنوان گزارش} \\
			[0.4cm]
			
			\large \bfseries
			نام درس \\ 
			[0.4cm]}
		\HRule \\[4cm]
		
		\large\textbf{نویسنده: }\\[0.2cm]
		نام و نام خوانوادگی 
		\\[3cm]
		
		\vfill
		
		{\large \today}
		
	\end{center}
\end{titlepage}

% متن مربوط به گزارش را در این قسمت قرار دهید.
\twocolumn
\ptext[0-10]


% در صورت نیاز به مراجع و پیوست از دو خط و فایل های مربوطه استفاده فرمایید.
% پیوست خود را در این قسمت وارد کنید.
\appendix
\onecolumn
% شروع محیط پیوست 
\section{پیوست}
\ptext[10-20]

% مراجع خود را در این قسمت وارد کنید.
% دستوری برای کوچک کردن اندازه فونت‌ها 
\small
% شروع محیط مراجع 

\renewcommand\bibname{مراجع}
\setLTRbibitems
% برای ارجاع دهی خودکار محتوای فایل bibtex مراجع را در فایل Refrence/myreferences.bib قرار داده و ذخیره کنید.  و یا میتوانید دو خط زیر را کامنت کرده و از دستورات بعدی ان بزرای ارجاع دهی استفاده کنید.
\bibliographystyle{IEEEtran}
\bibliography{Refrence/myreferences.bib}
%\begin{thebibliography}{99}
%
%\resetlatinfont
%% تذکر: چنانچه از نسخه‌های قدیمی زی‌پرشین و بی‌دی استفاده می‌کنید، هنگام استفاده از سه دستور بالا، با خطا 
%% مواجه می‌شوید. بنابراین باید زی‌پرشین و بی‌دی خود را آپدیت کنید.
%% چنانچه مرجع فارسی هم دارید باید یا از بسته Persian-bib استفاده کنید و یا راهنمای bidi را مطالعه کنید.
%
%\bibitem {TE} { \sc S. T. Aldrich, E.  Enochs, J. R. Garc\'{i}a Rozas, L. Oyonarte,} {\sl Covers and envelopes in Grothendieck categories: flat covers of complexes with applications,} J. Algebra {\bf 243} (2001), 615-630.
%
%\bibitem {AJL}{ \sc L. Alonso, A. Jerem\'{i}oas, J. Lipman,} {\sl  Duality and flat base change on formal schemes,}
%in “Studies in Duality on Noetherian Formal Schemes and Non-Noetherian Ordinary
%Schemes,” Contemp. Math. {\bf 244}, Amer. Math. Soc., 1999, 3-90.
%
%\bibitem {AF} {\sc F. W. Anderson, K. R. Fuller,} {\sl Rings and categories of modules,} Graduate Texts in Mathematics 13, Springer-Verlag, Berlin-Heidelberg-New York, 1974.
%
%\bibitem {AEHS} {\sc J. Asadollahi, H. Eshraghi, R. Hafezi, Sh. Salarian,} {\sl On the homotopy categories of projective and injective representations of quivers,} J. Algebra {\bf 346} (2011), 101-115.
%
%
%
%
%\end{thebibliography}
%
%

\end{document}
% پیوست ها 
\documentclass[a4paper 11px twocolumn]{article}
% تنظیمات و پکیج های مورد نیاز در این فایل قرار دارد. برای تنظیم و افزودن پکیج به این فایل مراجعه کنید.
\usepackage{amsmath, amssymb, amsthm} % بسته‌های ریاضی و قضایای امسال
\usepackage[all]{xy} % برای ترسیم نمودارها
\usepackage[left=1.3cm, right=1.3cm, top=1cm, bottom=1.3cm, nohead, hmargin=2cm,vmargin=2.5cm]{geometry} % تنظیمات حاشیه‌ها
\usepackage{graphicx} % برای درج تصاویر
\usepackage{fancyhdr} % برای تنظیم سربرگ و پاورقی صفحات
\usepackage{fontspec} % برای انتخاب فونتها
\usepackage{framed} % برای ایجاد قاب برای متن
\usepackage{ } % برای شمارش تعداد صفحات
\usepackage{cite} % برای ارجاع به منابع
\usepackage{pdfpages} % برای درج صفحات PDF
\usepackage{ifthen} % برای شرطی کردن تنظیمات
\usepackage{xcolor, colortbl} % برای استفاده از رنگها
\usepackage{pifont} % برای استفاده از نمادها

\usepackage{enumitem} % برای سفارشی‌سازی محیط‌های شماره‌گذاری
\usepackage{stackengine} % برای ترتیب متون و اشیاء
\usepackage{tabularx} % برای جداول با عرض متغیر
\usepackage{multirow} % برای ترکیب سلول‌های چند ردیفی در جدول
\usepackage{fourier} % برای فونتهای تایمز
\usepackage{array} % برای سفارشی‌سازی جداول
\usepackage{makecell} % برای سفارشی‌سازی سلول‌ها در جدول
\usepackage{ptext} % تولید متن تصادفی
\usepackage{hyperref} % برای ایجاد لینکهای رنگی
\usepackage[nottoc]{tocbibind} % برای اضافه کردن فهرست مطالب به فهرست مطالب
\usepackage{bidipoem} % برای شعر فارسی
\usepackage{tikz} % برای ترسیم نمودارهای تعاملی
\usepackage{atbegshi} % برای کنترل تصاویر و نمودارها
\usetikzlibrary{calc} % کتابخانه‌ای برای ترسیم نمودارهای TikZ


\usepackage{fancyhdr} % Custom headers and footers


% چنانچه قصد پرینت گرفتن نوشته خود را دارید، خط بالا را غیرفعال و  از دستور زیر استفاده کنید چون در صورت استفاده از دستور زیر‌‌، 
% لینک‌ها به رنگ سیاه ظاهر خواهند شد که برای پرینت گرفتن، مناسب‌تر است
%\usepackage[pagebackref=false]{hyperref}


% بسته‌ لازم برای تنظیم سربرگ‌ها
% بسته‌ای برای ظاهر شدن «مراجع» و «نمایه» در فهرست مطالب

% دستورات مربوط به ایجاد نمایه
%\usepackage{makeidx}
%\makeindex
%%%%%%%%%%%%%%%%%%%%%%%%%%
% فراخوانی بسته زی‌پرشین و تعریف قلم فارسی و انگلیسی
\usepackage{xepersian}
\settextfont[Scale=1]{XB Niloofar}
%%%%%%%%%%%%%%%%%%%%%%%%%%
% چنانچه می‌خواهید اعداد در فرمول‌ها، انگلیسی باشد، خط زیر را غیرفعال کنید
\setdigitfont[Scale=1.2]{XB Niloofar}

\defpersianfont\nas[Scale=2]{IranNastaliq}

\defpersianfont\titr[Scale=1.5]{B Titr Bold}

%%%%%%%%%%%%%%%%%%%%%%%%%%


%%%%%%%%%%%%%%%%%%%%%%%%%%
% تعریف و نحوه ظاهر شدن عنوان قضیه‌ها، تعریف‌ها، مثال‌ها و ...
\theoremstyle{definition}
\newtheorem{definition}{تعریف}[section]
\theoremstyle{theorem}
%\newtheorem{s}[definition]{}
\newtheorem{theorem}[definition]{قضیه}
\newtheorem{lemma}[definition]{لم}
\newtheorem{proposition}[definition]{گزاره}
\newtheorem{corollary}[definition]{نتیجه}
\newtheorem{remark}[definition]{تبصره}
\theoremstyle{definition}
\newtheorem{example}[definition]{مثال}
\newtheorem{notation}[definition]{نمادگذاری}
\newtheorem{definitiont}[definition]{تعریف و نمادگذاری}
\newtheorem{obs}[definition]{ملاحظه}
%%%%%%%%%%%%%%%%%%%%%%%%%%%%

\begin{document}
% مشخصات گزارش را در این قسمت تغییر بدهید.
\begin{titlepage}
	\setlength{\arrayrulewidth}{0.4mm}
	\newcommand{\HRule}{\rule{\linewidth}{0.5mm}}
	\begin{center}
		% logo
		\includegraphics[width=0.5\textwidth]{./Fig/logo.png}~\\[4cm]
		
		\HRule \\[0.4cm]
		{
			\Huge\textbf{عنوان گزارش} \\
			[0.4cm]
			
			\large \bfseries
			نام درس \\ 
			[0.4cm]}
		\HRule \\[4cm]
		
		\large\textbf{نویسنده: }\\[0.2cm]
		نام و نام خوانوادگی 
		\\[3cm]
		
		\vfill
		
		{\large \today}
		
	\end{center}
\end{titlepage}

% متن مربوط به گزارش را در این قسمت قرار دهید.
\twocolumn
\ptext[0-10]


% در صورت نیاز به مراجع و پیوست از دو خط و فایل های مربوطه استفاده فرمایید.
% پیوست خود را در این قسمت وارد کنید.
\appendix
\onecolumn
% شروع محیط پیوست 
\section{پیوست}
\ptext[10-20]

% مراجع خود را در این قسمت وارد کنید.
% دستوری برای کوچک کردن اندازه فونت‌ها 
\small
% شروع محیط مراجع 

\renewcommand\bibname{مراجع}
\setLTRbibitems
% برای ارجاع دهی خودکار محتوای فایل bibtex مراجع را در فایل Refrence/myreferences.bib قرار داده و ذخیره کنید.  و یا میتوانید دو خط زیر را کامنت کرده و از دستورات بعدی ان بزرای ارجاع دهی استفاده کنید.
\bibliographystyle{IEEEtran}
\bibliography{Refrence/myreferences.bib}
%\begin{thebibliography}{99}
%
%\resetlatinfont
%% تذکر: چنانچه از نسخه‌های قدیمی زی‌پرشین و بی‌دی استفاده می‌کنید، هنگام استفاده از سه دستور بالا، با خطا 
%% مواجه می‌شوید. بنابراین باید زی‌پرشین و بی‌دی خود را آپدیت کنید.
%% چنانچه مرجع فارسی هم دارید باید یا از بسته Persian-bib استفاده کنید و یا راهنمای bidi را مطالعه کنید.
%
%\bibitem {TE} { \sc S. T. Aldrich, E.  Enochs, J. R. Garc\'{i}a Rozas, L. Oyonarte,} {\sl Covers and envelopes in Grothendieck categories: flat covers of complexes with applications,} J. Algebra {\bf 243} (2001), 615-630.
%
%\bibitem {AJL}{ \sc L. Alonso, A. Jerem\'{i}oas, J. Lipman,} {\sl  Duality and flat base change on formal schemes,}
%in “Studies in Duality on Noetherian Formal Schemes and Non-Noetherian Ordinary
%Schemes,” Contemp. Math. {\bf 244}, Amer. Math. Soc., 1999, 3-90.
%
%\bibitem {AF} {\sc F. W. Anderson, K. R. Fuller,} {\sl Rings and categories of modules,} Graduate Texts in Mathematics 13, Springer-Verlag, Berlin-Heidelberg-New York, 1974.
%
%\bibitem {AEHS} {\sc J. Asadollahi, H. Eshraghi, R. Hafezi, Sh. Salarian,} {\sl On the homotopy categories of projective and injective representations of quivers,} J. Algebra {\bf 346} (2011), 101-115.
%
%
%
%
%\end{thebibliography}
%
%

\end{document}
% مراجع خود را در این قسمت وارد کنید.
% دستوری برای کوچک کردن اندازه فونت‌ها 
\small
% شروع محیط مراجع 

\renewcommand\bibname{مراجع}
\setLTRbibitems
% برای ارجاع دهی خودکار محتوای فایل bibtex مراجع را در فایل Refrence/myreferences.bib قرار داده و ذخیره کنید.  و یا میتوانید دو خط زیر را کامنت کرده و از دستورات بعدی ان بزرای ارجاع دهی استفاده کنید.
\bibliographystyle{IEEEtran}
\bibliography{Refrence/myreferences.bib}
%\begin{thebibliography}{99}
%
%\resetlatinfont
%% تذکر: چنانچه از نسخه‌های قدیمی زی‌پرشین و بی‌دی استفاده می‌کنید، هنگام استفاده از سه دستور بالا، با خطا 
%% مواجه می‌شوید. بنابراین باید زی‌پرشین و بی‌دی خود را آپدیت کنید.
%% چنانچه مرجع فارسی هم دارید باید یا از بسته Persian-bib استفاده کنید و یا راهنمای bidi را مطالعه کنید.
%
%\bibitem {TE} { \sc S. T. Aldrich, E.  Enochs, J. R. Garc\'{i}a Rozas, L. Oyonarte,} {\sl Covers and envelopes in Grothendieck categories: flat covers of complexes with applications,} J. Algebra {\bf 243} (2001), 615-630.
%
%\bibitem {AJL}{ \sc L. Alonso, A. Jerem\'{i}oas, J. Lipman,} {\sl  Duality and flat base change on formal schemes,}
%in “Studies in Duality on Noetherian Formal Schemes and Non-Noetherian Ordinary
%Schemes,” Contemp. Math. {\bf 244}, Amer. Math. Soc., 1999, 3-90.
%
%\bibitem {AF} {\sc F. W. Anderson, K. R. Fuller,} {\sl Rings and categories of modules,} Graduate Texts in Mathematics 13, Springer-Verlag, Berlin-Heidelberg-New York, 1974.
%
%\bibitem {AEHS} {\sc J. Asadollahi, H. Eshraghi, R. Hafezi, Sh. Salarian,} {\sl On the homotopy categories of projective and injective representations of quivers,} J. Algebra {\bf 346} (2011), 101-115.
%
%
%
%
%\end{thebibliography}
%
%

\baselineskip=1cm
\chapter*{واژه‌نامه فارسی به انگلیسی}\markboth{واژه‌نامه فارسی به انگلیسی}{واژه‌نامه فارسی به انگلیسی}
\addcontentsline{toc}{chapter}{واژه‌نامه فارسی به انگلیسی}

\englishgloss{Hyperparameter}{ابرپارامترها}
\chapter*{واژه‌نامه  انگلیسی به  فارسی}\markboth{واژه‌نامه  انگلیسی به  فارسی}{واژه‌نامه  انگلیسی به  فارسی}
\addcontentsline{toc}{chapter}{واژه‌نامه  انگلیسی به  فارسی}
\thispagestyle{empty}

\persiangloss{عامل}{Agent}

\printindex
\chapter*{فهرست اسامی}\markboth{فهرست اسامی}{فهرست اسامی }

\thispagestyle{empty}

\persiangloss{اوسلندر}{Auslander}


\baselineskip=.750cm
\thispagestyle{empty}
\begin{latin}

~~
\vspace{1.5cm}

\begin{flushleft}
\begin{Huge}
{\bf Abstract}
\end{Huge}
\end{flushleft} 
\vspace{0.8cm}
This thesis is devoted to study the category of representations of a guiver. 
\vspace{0.75cm}
\\{\bf Key Words.} Auslander-Reiren theory, almost split sequences for complexes,...
\end{latin}
% در این فایل، عنوان پایان‌نامه، مشخصات خود و چکیده پایان‌نامه را به انگلیسی، وارد کنید.
% توجه داشته باشید که جدول حاوی مشخصات پروژه/پایان‌نامه/رساله، به طور خودکار، رسم می‌شود.
%%%%%%%%%%%%%%%%%%%%%%%%%%%%%%%%%%%%
\begin{titlepage}
\begin{latin}
\begin{center}
\vspace{-2cm}
\centerline{{\includegraphics[height=4cm]{./BaseFile/logo.jpg}}}

\vspace{-1.2cm}
\textbf{\Timmm{University of Isfahan }}

\vspace{-.5cm}
\textbf{\Timm{ Faculty of Mathematics and Statistics}}

\vspace{-.5cm}
\textbf{\Tim{Department of Applied Mathematics and Computer Science}}

\vspace{1.2cm}
\textbf{\Timmmmmm{MSc. Thesis }}

\vspace{.5cm}
\textbf{\Timmmmmm{ Title }}


\vspace{1cm}
\textbf{\Timmmmm{Supervisor:}}

\textbf{\Timmmm{Dr. ...}}

\textbf{\Timmmm{Dr. ...}}
\vspace{.6cm}

\textbf{\Timmmmm{Advisor:}}

\textbf{\Timmmm{Dr. ....}}

\vspace{.6cm}
\textbf{\Timmmmm{By:}}

\textbf{\Timmmm{ \fullnameEN }}


\vspace{.6cm}
\textbf{\Timmmm June  2025}
\end{center}
\end{latin}
\end{titlepage}

\label{LastPage}
\end{document}

